
lieu de la th\'eorie de Kummer le fait suivant, d\^u \`a Grothendieck (SGA 1 XI, \`a la derni\`ere page)~: si @@M0@@ est un sch\'ema propre sur un corps @@M1@@, avec @@M2@@, alors il existe un isomorphisme de @@M3@@-sch\'emas en groupes @@M4@@; par cons\'equent, si @@M5@@ est alg\'ebriquement clos, on a @@M6@@.
\item Si @@M7@@ est de caract\'eristique nulle, alors le noyau @@M8@@ est nul, puisque tout groupe fini unipotent est alors nul (SGA 3 XVII 1.7). On retrouve ainsi \`a peu pr\`es le ``Vanishing theorem'' de Kodaira sous la forme donn\'ee par Mumford [2].
\end{enumerate}

\textbf{Th\'eor\`eme 3.13.} \emph{Soient @@M9@@ un sch\'ema noeth\'erien et @@M10@@ un @@M11@@-sch\'ema projectif, muni d'un faisceau inversible relativement ample @@M12@@. On suppose que les fibres g\'eom\'etriques de @@M13@@ soient int\`egres et de dimension @@M14@@. Alors pour une famille @@M15@@ de classes de faisceaux inversibles sur les fibres de @@M16@@, les conditions suivantes sont \'equivalentes}~:
\begin{enumerate}[label=(\roman*)]
\item \emph{@@M17@@ est une famille limit\'ee.}
\item \emph{Dans les polyn\^omes de Hilbert @@M18@@ des @@M19@@, le coefficient @@M20@@ reste born\'e et le coefficient @@M21@@ reste minor\'e.}
\item \emph{Lorsque @@M22@@ parcourt @@M23@@, le degr\'e @@M24@@ reste born\'e et le nombre d'intersection @@M25@@ reste minor\'e.}
\item \emph{Il existe un entier @@M26@@ tel que l'on ait}
@@M27@@
\emph{au sens de la relation d'ordre d\'efinie par la relation d'amplitude\,; c'est-\`a-dire, pour tout @@M28@@, @@M29@@ et @@M30@@ sont amples.}
\end{enumerate}

